% !TeX root = ../papers/nc-p3.tex
\section{可逆交差長方形の七つの台での正則性の決定}
\label{sec:parallel-support-fibers}
\begin{theorem}\label{thm:parallel-support-fibers}
平行台の代表82,118,124,129,134,139,142で正則になる窓は全て特別成分 $b=c=1,\ t=\pm1$ に属する。
各成分における正則点は次の表で尽きる。112型で $B=1,a=A/t,y=C,d=D/t$ とし、表の自由な文字には
指定した台の他の係数が非零という条件を課す。
\[
\begin{array}{c|l|l}
 台&\text{特別成分内の等式}&\text{正則点}\\\hline
82&A=B&\text{全点}\\
118&a=y=1,\ d=0&\text{全点}\\
124&a=y=1&\text{全点}\\
129&a=1,\ y=0&d\ne1\\
134&a=1,\ d=0&\text{全点}\\
139&a=1&d\ne1\\
142&d=0,\ a+y=1&\text{全点}
\end{array}
\]
台85,145,147の特別成分上の点も全て正則であり、台143,144,146は特別成分と交わらない。
\end{theorem}
\begin{proof}[全厚い係数の零条件]
特別成分の112型で16成分の零条件は順に
\[
 a,1,a,d,1,y,d,y,\ell,a-1,\ell,y-1,a-1,1,y-1,1,\quad
 \ell=a+y-d-1
\]
の消滅と同値で、$ay-d\ne0$ である。76型でも順に
\[
 A,B,A,D,B,C,D,C,C-D,A-B,C-D,C,A-B,B,C,B
\]
となる。各々、右下八成分の共通分母と単元符号を除いただけである。
従って定理\ref{thm:parallel-invertible-rectangles}の四つの例外直線対を各台に制限すると表の最終列を得る。
例えば112型の対象の台では $a\ne0$ なので障害は
$(a=d=1)$、$(y=d=1)$、$(y=0,a=d)$ の三条件に簡約される。

主成分では $B=1$ とする。76型では $C=1$ で、$x_3$ と $x_9$ は同時に、$x_6$ と $x_{12}$ も同時に消える。
従って $x_9,x_{12}$ だけが零の台82は主成分にない。
112型では $a=A/t,y=C,m=t^2=b+c-bc$ と置く。$n=-c(a-b)(y-b)$ とすれば、
各成分の零条件は次の16多項式の消滅とそれぞれ同値である。
\[
\begin{gathered}
 ca+b(1-c),\ 1,\ a,\ n+b^2,\ 1,\ y,\ c(a-1)(y-1)-m,\ cy+b(1-c),\\
 ay-b,\ a-b,\ n+b(1+b-a-y),\ y-1,\ a-1,\ 1,\ y-b,\ 1 .
\end{gathered}
\]
実際、各分子を $t^2=m$ で簡約すれば表示式との比は単元であり、分母には
$tbc$ と行列式因子 $ay(b+c)-bc(a+y)+b^2(c-1)$ だけを要する。

台118,124,129,134,139は $x_9=x_{12}=0$ を含むため、$a=b=1,t^2=1$ である。
すると $D=t$ で、左上の行列式条件は $y\ne1$ となる。これで118,124,134は主成分上で不可能である。
129では $y=0$ と $cy+1-c=0$ も要るので $c=1$ となる。
残る139については、より強く $a=b=1$ の二つの主成分切片の全点が不正則である。
実際、$t=\pm1$ それぞれで、$\mathbb Q(c,y)$ 上の始点1の長さ7の四文字の次元は
$(34,33,34,33)$、合計134である。標準値は120なので上半連続性によりあらゆる特殊値でも排除される。
計算は各次数の自由語基底に一文字を付加し、末尾での六関係の行空間を割る有理関数体上の掃き出しである。

台142では $x_8=x_{10}=0$ かつ $x_{11}x_{14}\ne0$ より、$ay=b$ および $c=1$ を得る。
実際、$a=b/y$ のとき $x_{10}$ を検出する多項式は $b(y-b)(y-1)(c-1)/y$ である。
$x_3=x_6=0$ よりさらに $a+y=1=b$ となり、特別成分上に戻る。
\end{proof}

\begin{corollary}
主成分で台143は $t=\pm1$, $c=1$, $a=1-y$, $b=y(1-y)$ の二曲線の開部分、
145は $t=\pm1$, $c=1$, $a=b/y$ の二曲面の開部分である。
144は $a+y=1$, $c=b^2/(ay-b+b^2)$ の切片、146は
$c=b(1+b-a-y)/((a-b)(y-b))$ の切片である。いずれも元の台の非零条件と $t^2=b+c-bc$ を課す。
両形の全パラメータで $t$ の符号変更は四元群の二コサイクル捻りと位置別対角変更で実現される。
112型では同時に $A,D$ の符号を変え、76型では $A,B,C,D$ を固定する。
\end{corollary}
\begin{proof}
零条件表で $x_{10}$ と $x_3$ の多項式の差は $-b(a+y-1)$ である。
これと上の $ay=b$ での式を使えば各切片の表示を得る。
捻りについては、符号変更前後の一般16成分と薄い成分の比が全て定数符号になる。
各語でこの比に四元群二コサイクルの符号を掛けると、位置別倍率の重み行列の全左整数核で指標値が1である。
従って一組の定数対角倍率で全パラメータに共通の同値を与え、零成分を持つ場合にも特殊化する。
\path{computations/verify_parallel_support_fibers.py} が係数の単元、台、両符号の長さ7と捻り指標を検算する。
\end{proof}
